% !TEX root = main.tex
%%%%%%%%%%%%%%%%%%%%%%%%%%%%%%%%%%%%%%%%%%%%%%%%%%%%%%%%%%%%%%%%%%%%%%%
% FRONT MATTER - title, abstract, contents
%%%%%%%%%%%%%%%%%%%%%%%%%%%%%%%%%%%%%%%%%%%%%%%%%%%%%%%%%%%%%%%%%%%%%%%

\title{\bfseries Endogenous Investment and Capital Accumulation\\[2pt] in a Disaggregated Keynesian Cross\\[6pt] \large A Normalized-CES Agent-Based Model with an Endogenous Labor Market}

\author{%
  Mattia \textsc{Saramin}\thanks{%
  Correspondence: \texttt{904394@stud.unive.it}. This paper documents an exploratory computational model. Every number reported here is either measured from a committed, reproducible driver script, or explicitly declared as a design target or convention; the distinction is maintained throughout (\cref{sec:method}). Code, data and figures: \url{https://github.com/MattiaSaramin/endogenous-investment-abm}.}\\[2pt] \normalsize Department of Economics, Ca' Foscari University of Venice}

\date{}

\maketitle


\begin{abstract}
\noindent
We extend \citet{Teglio2025}'s disaggregated Keynesian cross with endogenous investment and capital accumulation. Production is a normalized CES \citep{KlumpDeLaGrandville2000} with elasticity~$\sigma$; firms hire endogenously, the wage follows a Blanchflower-Oswald curve, and the model is stock-flow consistent throughout.

Three results: (1) $Y(\rho)$ is U-shaped: its turning point $\rhostar$ rises with $\sigma$, and the anchored retention rate lies left of that turn at every elasticity tested, so the marginal effect is negative, the wide-range effect positive. Through a single slope this is a sign frontier $\sigstar \approx 0.65$, a declared reduction, not a verdict. (2) The region where retention depresses output is conditional on fiscal institutions: a balanced-budget benefit at an OECD-band replacement rate removes it. (3) The wage curve generates a fragility immune to demand instruments: three stabilization hypotheses fail through one wage $\,\to\,$unemployment$\,\to\,$ capital-erosion channel. That immunity is diagnostic, the fragility is a relative-price phenomenon, not a demand one, and it is bracketed by the price convention: it stands under a fixed numéraire, and is switched off, through a Pigou effect, under complete instantaneous pass-through, with neither convention anchored.

A sixteen-parameter global sensitivity analysis separates what generalises from what does not. The negative margin at the anchored $\rho$ survives marginalization where the turn resolves. The frontier does not: $\sigma$ is nearly inert once other parameters move, so a sign crossing there is conditional, not structural.
Viability is governed almost entirely by the depreciation rate, at whose empirically implied value no configuration survives.

\bigskip
\noindent\textbf{Keywords:} agent-based macroeconomics; stock-flow consistency; elasticity of substitution; wage-led growth; global sensitivity
analysis; capital accumulation.

\smallskip
\noindent\textbf{JEL:} E11, E12, E22, C63, D31.
\end{abstract}

\clearpage
\tableofcontents
\clearpage
